\documentclass[a4paper,12pt]{article}
\usepackage[utf8]{inputenc}
\usepackage[russian]{babel}
\usepackage{amsmath, amssymb, amsthm}
\usepackage{geometry}
\geometry{left=2cm, right=2cm, top=2cm, bottom=2cm}

\title{}
\author{}
\date{}

\begin{document}
	
	\maketitle
	
	\section*{Основная часть}
	
	\subsection*{Задача 1.}
	
	\subsubsection*{1) $f(x)=\max\{|x|,x^2\}$, $x\in\mathbb{R}$}
	
	Заметим, что
	\[
	\max\{|x|,x^2\}=\begin{cases}
		|x|, & |x|\le1,\\[1mm]
		x^2, & |x|\ge1.
	\end{cases}
	\]
	При поиске супремума разбиваем задачу на два региона.
	
	\textbf{Регион 1: $|x|\le1$.}  
	При $x\ge0$ имеем
	\[
	xy - f(x) = xy - x = x(y-1).
	\]
	Максимум по $x\in[0,1]$ равен:
	\[
	\begin{cases}
		y-1, & \text{если } y\ge1,\\[1mm]
		0, & \text{если } y<1.
	\end{cases}
	\]
	При $x<0$, поскольку $|x|=-x$, получаем:
	\[
	xy - f(x) = xy - (-x)= x(y+1).
	\]
	Максимум по $x\in[-1,0]$ равен:
	\[
	\begin{cases}
		-y-1, & \text{если } y\le -1,\\[1mm]
		0, & \text{если } y>-1.
	\end{cases}
	\]
	
	\textbf{Регион 2: $|x|\ge1$.}  
	Здесь $f(x)=x^2$ и функция
	\[
	\varphi(x)=xy-x^2
	\]
	есть квадратичная, и её оптимум по $\mathbb{R}$ достигается при $x=\frac{y}{2}$. При этом условие $|x|\ge1$ выполнится, если $|y|\ge2$. Тогда
	\[
	\varphi\left(\frac{y}{2}\right)=\frac{y^2}{2}-\frac{y^2}{4}=\frac{y^2}{4}.
	\]
	Если $|y|<2$, оптимум достигается на границе: при $y\ge0$ на $x=1$ даётся значение $y-1$, а при $y<0$ на $x=-1$ --- значение $-y-1$.
	
	\textbf{Итог.} Сопоставляя оба региона, получаем окончательный ответ:
	\[
	f^*(y)=
	\begin{cases}
		\frac{y^2}{4}, & |y|\ge2,\\[1mm]
		y-1, & 1\le y<2,\\[1mm]
		0, & |y|<1,\\[1mm]
		-y-1, & -2<y\le -1.
	\end{cases}
	\]
	При проверке на стыках видно, что функция непрерывна.
	
	\vspace{2mm}
	
	\subsubsection*{2) $f(x)=(|x|+1)\ln(|x|+1)$, $x\in\mathbb{R}$}
	
	Заметим, что функция чётная, поэтому можно положить $t=|x|\ge0$. Тогда
	\[
	f(x)= (t+1)\ln(t+1),
	\]
	а сопряжённая функция определяется как
	\[
	f^*(y)=\sup_{t\ge0}\Big\{t|y|-(t+1)\ln(t+1)\Big\}.
	\]
	Обозначим
	\[
	\phi(t)= t|y|-(t+1)\ln(t+1).
	\]
	Находим оптимальную точку. Вычислим производную:
	\[
	\phi'(t)=|y|-\ln(t+1)-1.
	\]
	Приравнивая к нулю,
	\[
	\ln(t+1)=|y|-1 \quad \Longrightarrow \quad t+1=\exp(|y|-1),\quad t=\exp(|y|-1)-1.
	\]
	Заметим, что для $|y|<1$ при $t=0$ получаем $\phi(0)= -\ln1=0$, а производная в нуле равна $\phi'(0)=|y|-1<0$, поэтому максимум достигается при $t=0$ и равен 0.
	
	Таким образом, окончательный результат:
	\[
	f^*(y)=
	\begin{cases}
		0, & |y|<1,\\[1mm]
		\exp(|y|-1)-|y|, & |y|\ge1.
	\end{cases}
	\]
	
	\subsubsection*{3) $f(x)=(\|x\|+1)\ln(\|x\|+1)-\|x\|$, $x\in\mathbb{R}^d$}
	
	Так как функция зависит только от нормы, оптимизация по $x\in\mathbb{R}^d$ сводится к одномерной задаче по $t=\|x\|\ge0$. При выборе направления оптимально брать $x$ коллинеарным с $y$, тогда
	\[
	\langle x,y\rangle = t\|y\|.
	\]
	Определяем функцию
	\[
	\psi(t)= t\|y\| - \Big[(t+1)\ln(t+1)-t\Big] = t(\|y\|+1) - (t+1)\ln(t+1).
	\]
	Найдём оптимальную точку:
	\[
	\psi'(t)=\|y\|+1 - \big[\ln(t+1)+1\big] = \|y\|-\ln(t+1).
	\]
	Приравнивая к нулю, получаем:
	\[
	\ln(t+1)=\|y\| \quad \Longrightarrow \quad t+1=\exp(\|y\|),\quad t=\exp(\|y\|)-1.
	\]
	Подставляем в $\psi(t)$:
	\[
	\psi\Big(\exp(\|y\|)-1\Big)= (\exp(\|y\|)-1)(\|y\|+1) - \exp(\|y\|)\|y\| 
	=\exp(\|y\|)-\|y\|-1.
	\]
	Таким образом, сопряжённая функция имеет вид:
	\[
	f^*(y)=\exp(\|y\|)-\|y\|-1.
	\]
	
	\vspace{3mm}
	
	\subsection*{Задача 2.}
	
	Пусть
	\[
	f(x)=\cosh x -1 = \frac{e^x+e^{-x}}{2}-1,\quad x\in\mathbb{R}.
	\]
	Найдём сопряжённую функцию
	\[
	f^*(y)=\sup_{x\in\mathbb{R}}\Big\{xy - (\cosh x-1)\Big\}.
	\]
	Так как $f$ строго выпуклая и дифференцируемая, оптимальная точка определяется условием:
	\[
	\frac{d}{dx}\Bigl(xy-\cosh x+1\Bigr)= y-\sinh x=0\quad \Longrightarrow\quad x=\operatorname{arcsinh}(y).
	\]
	Подставляя, получаем:
	\[
	f^*(y)=y\operatorname{arcsinh}(y)-\cosh(\operatorname{arcsinh}(y))+1.
	\]
	Используем известное равенство
	\[
	\cosh(\operatorname{arcsinh}(y))=\sqrt{1+y^2},
	\]
	и окончательно получаем:
	\[
	f^*(y)=y\operatorname{arcsinh}(y)-\sqrt{1+y^2}+1.
	\]
	Заметим, что так как $f$ --- замкнутая выпуклая функция, то её двойственное сопряжение совпадает с $f$, то есть
	\[
	f^{**}(x)=f(x)=\cosh x-1.
	\]
	
	\section*{Дополнительная часть}
	
	\subsection*{Задача 1.}
	
	Здесь принято, что функция
	\[
	f(x)=\sum_{i=1}^d\log x_i
	\]
	определена на множестве $x=(x_1,\dots,x_d)$, где $x_i>0$, $i=1,\dots,d$. Тогда сопряжённая функция определяется как
	\[
	f^*(y)=\sup_{x>0}\left\{\langle y,x\rangle-\sum_{i=1}^d\log x_i\right\}.
	\]
	Благодаря аддитивности оптимизация разбивается по координатам. Для каждой координаты рассматриваем задачу
	\[
	\sup_{x>0}\{y_i x-\log x\}.
	\]
	Для фиксированного $i$ оптимум находится по условию первого порядка:
	\[
	\frac{d}{dx}\Bigl(y_i x-\log x\Bigr)=y_i-\frac{1}{x}=0\quad\Longrightarrow\quad x=\frac{1}{y_i},
	\]
	при условии, что $y_i>0$. Подставляя обратно,
	\[
	y_i\cdot \frac{1}{y_i}-\log\frac{1}{y_i}=1+\log y_i.
	\]
	Если хотя бы для одного $i$ условие $y_i>0$ не выполнено, супремум равен $+\infty$. Таким образом,
	\[
	f^*(y)=
	\begin{cases}
		\displaystyle \sum_{i=1}^d\Bigl(1+\log y_i\Bigr), & \text{если } y_i>0\ \forall i,\\[2mm]
		+\infty, & \text{иначе.}
	\end{cases}
	\]
	
	\subsection*{Задача 2.}
	
	\medskip
	
	\textbf{1) Если $f(x)=\alpha\, g(x)$, $\alpha>0$,} то
	\[
	f^*(y)=\sup_{x}\{ \langle y,x\rangle-\alpha\, g(x)\}
	=\alpha\,\sup_{x}\Bigl\{\Bigl\langle \frac{y}{\alpha},x\Bigr\rangle-g(x)\Bigr\}
	=\alpha\, g^*\Bigl(\frac{y}{\alpha}\Bigr).
	\]
	
	\medskip
	
	\textbf{2) Если $f(x)=g(Ax)$,} где $A\in S_d^{++}$ (матрица симметричная, положительно определённая), делаем замену $u=Ax$. Тогда $x=A^{-1}u$, и
	\[
	\langle y,x\rangle=\langle y, A^{-1}u\rangle=\langle A^{-1}y,u\rangle.
	\]
	Следовательно,
	\[
	f^*(y)=\sup_{u}\Bigl\{\langle A^{-1}y,u\rangle-g(u)\Bigr\}=g^*(A^{-1}y).
	\]
	
	\medskip
	
	\textbf{3) Если $f(x)=\inf\{g(u)+h(v):\, u+v=x\}$,} то по известному результату сопряжение инфимум-суммы даёт
	\[
	f^*(y)=g^*(y)+h^*(y).
	\]
	
	\bigskip

\end{document}
